% contributions.tex -- research questions and contributions.
% 4 Oct draft; restructured 9 Oct (structure audit); 10 Oct Wesley's five questions.
% 11 Oct: the adopted spine (docs/thesis_plan_2026-10-11.md, Wesley 13:40): RQ1 carry, RQ2 leave, RQ3 join;
% the prediction ledger is Section 4.4, not a question; the planner-against-network question of the project
% plan is a scope paragraph plus Appendix F. Contributions C1-C3 and method items M1-M2 from the plan's
% section 2.4. Numbers that are still moving (the join claim table, flown Fri 16; the rig attach, 21-22 Oct)
% are left as bracketed placeholders, per the plan: prose starts Fri 16.

\subsection{Research Questions}
\label{sec:rq}

The research question of this thesis is whether a centralised controller can let a quadcopter leave or join
a team that is carrying a payload in flight, and what sets the payload's transient when it does. The payload
is a 0.86\,kg ring carried on rigid rods (Section~\ref{sec:hardware}). The controller is a cascade: a
centralised planner solves an optimal control problem (OCP) for the payload and every rod, and a tracking
controller per UAV follows its share of the plan. A UAV leaving or joining resizes the planner's OCP to the
new team, a switch to a pre-built problem with the attachment geometry as a runtime parameter
(Section~\ref{sec:reconfig}); a joining UAV is called the newcomer. One question establishes the carry, two
test the reconfiguration.
\begin{itemize}[leftmargin=3.6em, labelwidth=3.0em, labelsep=0.6em, align=left]
    \item[\textbf{RQ1}] \textbf{Carry.} What does feeding the trackers the rod pull measured from motion
    capture buy on unmodified Betaflight firmware, and which sources of steady payload error does the rig
    expose that the twin does not?
    \item[\textbf{RQ2}] \textbf{Leave.} When a carrier leaves in flight, what sets the ring's transient, and what
    keeps the ring from being lost?
    \item[\textbf{RQ3}] \textbf{Join.} What must be observed before a quadcopter flown by another controller is
    admitted to a carried ring, and what does admission cost the ring?
\end{itemize}
RQ1 is the baseline for the other two: a reconfiguration can only be judged against a team that carries the
payload well before anything changes. The working hypothesis is that the centralised approach suits a change
of team composition because the coordinating node solves for the force distribution of the team it actually
has in the first tick after the event, instead of leaving the survivors to find it from the motion that
follows. Beside the three questions, Section~\ref{sec:ledger} sets up a ledger that will score the rig twin's
prediction for each rig cell that gets a band frozen before it flies (none is scored yet).

The project plan's secondary question, whether a decentralised network controller of the kind of Quan et
al.\ \cite{quan2025swarm} or the trajectory-based OCP \cite{tudelft} better balances robustness and tracking, was
answered on an earlier simulation plant, before the rig twin existed, and is not re-flown:
Appendix~\ref{sec:app_network} records it as dated history.

Table~\ref{tab:rq_map} gives the experiments behind each question, flown in the rig twin (a Gazebo simulation
of the lab rig, Section~\ref{sec:simulation}) and on the rig, and the section that answers it.

\begin{table}[h]
\centering
\small
\caption{The research questions, the experiments behind them and the sections that answer them.}
\label{tab:rq_map}
\begin{tabularx}{\textwidth}{>{\raggedright\arraybackslash}p{2.6cm}>{\raggedright\arraybackslash}Xl}
\toprule
Question & Experiments & Answered in \\
\midrule
RQ1 Carry & In the twin, three UAVs from a floor start along five load paths at three speeds, and with 1\,m
cables; pushes, wrong model values and the tracker's initial condition. On the rig, four UAVs on the circle and
the figure-eight with the planned and the measured rod pull; the thrust map; the steady errors the rig showed &
Section~\ref{sec:res_carry} \\
\addlinespace
RQ2 Leave & Every release on the rig: announced, unannounced, and the magnets that let go with no command; in the
twin, the levers on the transient: the survivors' pull source, detecting the release, the support margin; the
switch itself; the drop rule & Section~\ref{sec:res_leave} \\
\addlinespace
RQ3 Join & The admission contract with the collaborator's controller: contact, the collaborator's lift,
confirmation from motion capture alone, the hand-over and the fold-in; the take-back when nothing is proven; the
attach on the rig & Section~\ref{sec:res_join} \\
\addlinespace
The prediction ledger & Each rig cell against the twin's band for it, frozen before the flight (none scored yet) &
Section~\ref{sec:ledger} \\
\addlinespace
All three against the literature & -- & Section~\ref{sec:discussion} \\
\bottomrule
\end{tabularx}
\end{table}

\subsection{Contributions}
\label{sec:contributions}

Three claims and two method items, in the order of the questions.
\begin{enumerate}
    \item \textbf{C1. Carrying with a measured rod pull on unmodified Betaflight firmware (RQ1).} The
    load--cable planner of \cite{tudelft} is combined with a cable-aware tracking model predictive controller
    (MPC) per UAV that commands Betaflight flight controllers over their radio links, with one affine thrust map
    identified on the rig for all four airframes, a floor start for rigid rods, and each tracker fed the rod pull
    it measures from motion capture (Sections~\ref{sec:controller}, \ref{sec:tracker} and~\ref{sec:takeoff}). On
    the rig four UAVs carried the 0.86\,kg ring around a circle and through a figure-eight within 3\,cm of the
    path on average once the trackers used the measured pull and the rods were planned at 60$^\circ$, against
    about 8\,cm before. In the rig twin three UAVs carried it along five load paths at peak path speeds up to
    0.3\,m/s inside bars frozen before the runs in 39 of 40 runs, and in this rigid-rod twin the measured pull took up a 0.6\,N
    vertical push on the ring within about a second, with no integral action (Section~\ref{sec:res_carry}).
    \item \textbf{C2. The leave transient on hardware, and the levers that set it (RQ2).} Six commanded
    releases on the rig: the three remaining UAVs carried the ring on through five, with tilt rises of
    5.3--16.9$^\circ$, and unannounced releases were found from the stretched rod in 0.33--0.42\,s. Every ring
    lost on the rig began with a magnet letting go. In the twin the tilt rise grew as the support margin the leaver
    leaves behind shrank, and in two hover cells it was smaller with the survivors on the measured pull than on the
    planner's pull (Section~\ref{sec:res_leave}). A geometric condition, checkable
    before a flight, says which carrier a layout may lose: a four-carrier layout admits at most two leavers and
    never all four, so a team that must tolerate any single loss needs five (Section~\ref{sec:reconfig}).
    \item \textbf{C3. An admission contract for a UAV flown by another controller (RQ3).} The other
    controller descends to contact, lifts and holds; this work confirms the catch from motion capture alone, from
    the magnet tip moving with the plate on a taut tether, switches the radio to its own tracker on the same tick,
    folds the newcomer in from a small starting tension, and takes it back and lands it if nothing is proven
    (Section~\ref{sec:admission}). [Twin claim table: the join code tree, after the 21--22 October rig day. Rig:
    eight attach flights on 14--15 October gave one contact and no fold-in; the attach on 21--22 October.]
    (Section~\ref{sec:res_join}).
    \item \textbf{M1. Reconfiguration and authority.} Notice, re-plan, redistribute, fall back: the attachment
    geometry as a runtime parameter with one pre-built solver per team size, primed before it flies; exactly
    one owner per UAV at every instant, and a way down from every phase, each flown in the twin (the attach
    phases on the earlier attach sequence): the take-back, the tilt cut,
    LAND for every UAV, landing the fleet on a newcomer fault, dropping the ring on a second loss or a tilt
    over 30$^\circ$ (which fired on the rig after a magnet let go), and the freed UAV disarmed alone
    (Sections~\ref{sec:reconfig} and~\ref{sec:safety}).
    \item \textbf{M2. The claim protocol and the prediction ledger.} Every claim is flown against a named
    baseline with its bars frozen before the runs, on one code tree, and each rig cell gets a twin band frozen
    before it flies, to be scored afterwards (no rig row is scored yet); where the twin and the rig have been
    compared, it matched the reactive release and missed what it does not model (the magnet's hold, rotor wake,
    the rig's latency)
    (Sections~\ref{sec:experiments} and~\ref{sec:ledger}).
\end{enumerate}

\subsection{Outline}

Section~\ref{sec:lit} reviews the literature and states the gap. Section~\ref{sec:method} describes the
rig and its twin, the model, the controller, the floor start and landing, the reconfiguration and its authority
rules, and the experimental method. Section~\ref{sec:results} answers the questions in the order of
Table~\ref{tab:rq_map}, keeps the twin's ledger and discusses the answers against the literature, and
Section~\ref{sec:conclusion} concludes. The appendices hold the software, the parameters, the rig operations,
the secondary and negative results, the development timeline, and the network controller of the earlier plant.
